% LuaHBTeX 1.14 で生成したPDFを古いPopplerでも安定して表示できるようにします。
\pdfvariable objcompresslevel=0

% 日本語と書体
\usepackage{luatexja-fontspec}
\usepackage{fontspec}
\setmainfont{Noto Serif CJK JP}[AutoFakeSlant=0.2]
\setsansfont{Noto Sans CJK JP}[AutoFakeSlant=0.2]
\setmonofont{DejaVu Sans Mono}[Scale=0.86]
\setmainjfont{Noto Serif CJK JP}
\setsansjfont{Noto Sans CJK JP}
\setmonojfont{Noto Sans Mono CJK JP}[Scale=0.86]

% ページと本文
\usepackage[
  margin=25mm,
  headheight=15pt,
  footskip=12pt
]{geometry}
\usepackage{setspace}
\usepackage{indentfirst}
\setstretch{1.12}
\setlength{\parindent}{1\zw}
\setlength{\parskip}{0.22\baselineskip}
\setlength{\emergencystretch}{3em}
\raggedbottom

% 数式、表、箇条書き
\usepackage{fix-cm}
\usepackage{amsmath,amssymb,mathtools}
\usepackage{booktabs,longtable,tabularx,array}
\usepackage{enumitem}
\usepackage{makeidx}
\makeindex
\renewcommand{\indexname}{索引}
\setlist{leftmargin=2.2em,itemsep=0.25em,topsep=0.45em}
\newcolumntype{Y}{>{\raggedright\arraybackslash}X}

% 図
\usepackage{graphicx}
\usepackage{float}
\usepackage{placeins}
\usepackage{tikz}
\usetikzlibrary{arrows.meta,automata,calc,fit,matrix,positioning,shapes.geometric}

% 配色
\usepackage{xcolor}
\definecolor{OffWhite}{HTML}{F6F3F6}
\definecolor{DeepNavy}{HTML}{2C2437}
\definecolor{MutedBlue}{HTML}{5C4D7A}
\definecolor{PaleGray}{HTML}{71618E}
\definecolor{PaleBlue}{HTML}{E8CDDD}
\definecolor{CardWhite}{HTML}{FFFFFF}
\definecolor{RuleGray}{HTML}{D8CCD7}

% コード
\usepackage{listings}
\lstdefinestyle{codebase}{
  basicstyle=\ttfamily\small,
  keywordstyle=\color{MutedBlue}\bfseries,
  commentstyle=\color{PaleGray},
  stringstyle=\color{DeepNavy},
  showstringspaces=false,
  columns=fullflexible,
  keepspaces=true,
  breaklines=true,
  breakatwhitespace=false,
  frame=single,
  rulecolor=\color{RuleGray},
  backgroundcolor=\color{OffWhite},
  xleftmargin=2mm,
  xrightmargin=2mm,
  aboveskip=0.8em,
  belowskip=0.8em
}
\lstdefinelanguage{Go}{
  morekeywords={package,import,func,var,const,type,struct,interface,map,range,return,if,else,for,switch,case,default,go,defer,select,chan},
  sensitive=true,
  morecomment=[l]{//},
  morecomment=[s]{/*}{*/},
  morestring=[b]"
}
\lstset{style=codebase}

% 注記欄
\usepackage[most]{tcolorbox}
\newtcolorbox{chapterabstractbox}{
  enhanced,
  breakable,
  colback=OffWhite,
  colframe=RuleGray,
  boxrule=0.5pt,
  arc=0pt,
  left=4mm,right=4mm,top=3mm,bottom=3mm,
  before skip=1.2em,after skip=1.2em
}
\newtcolorbox{takeawaybox}{
  enhanced,
  breakable,
  colback=PaleBlue,
  colframe=MutedBlue,
  boxrule=0pt,
  leftrule=2pt,
  arc=0pt,
  left=4mm,right=4mm,top=2.5mm,bottom=2.5mm,
  before skip=1em,after skip=1em
}
\newenvironment{learninggoals}{%
  \par\vspace{1.2em}{\Large\sffamily\bfseries\color{DeepNavy}この章の到達目標}\par\vspace{0.6em}
  \begin{itemize}
}{%
  \end{itemize}
}

% 見出しとヘッダー
\usepackage{titlesec}
\titleformat{\section}
  {\Large\sffamily\bfseries\color{DeepNavy}}
  {\thesection}{0.8em}{}
\titleformat{\subsection}
  {\large\sffamily\bfseries\color{DeepNavy}}
  {\thesubsection}{0.7em}{}
\titleformat{\paragraph}[runin]
  {\raggedright\normalsize\sffamily\bfseries\color{DeepNavy}}
  {}{0pt}{}[：]
\titlespacing*{\section}{0pt}{2.8em}{1.0em}
\titlespacing*{\subsection}{0pt}{2.0em}{0.7em}
\titlespacing*{\paragraph}{0pt}{1.0em}{0.5em}

% 図表の \centering を本文へ持ち越さず、行内見出しから通常の段落へ戻す
\let\BookParagraph\paragraph
\renewcommand{\paragraph}[1]{\par\raggedright\BookParagraph{#1}}

\usepackage{fancyhdr}
\pagestyle{fancy}
\fancyhf{}
\fancyhead[L]{\small\sffamily\nouppercase{\leftmark}}
\fancyhead[R]{\small\sffamily\thepage}
\renewcommand{\headrulewidth}{0.4pt}
\renewcommand{\footrulewidth}{0pt}
\fancypagestyle{plain}{%
  \fancyhf{}%
  \fancyhead[R]{\small\sffamily\thepage}%
  \renewcommand{\headrulewidth}{0pt}%
}
\renewcommand{\chaptermark}[1]{\markboth{第\thechapter 章\quad #1}{}}

% 目次
\setcounter{tocdepth}{1}
\setcounter{secnumdepth}{2}
\makeatletter
\renewcommand{\@pnumwidth}{2.5em}
\renewcommand{\@tocrmarg}{4em}
\makeatother

% リンクとPDFメタデータ
\usepackage[
  hidelinks,
  pdftitle={情報工学入門},
  pdfsubject={一つのAIアプリ要求を追い、物理からアプリまでをつなぐ技術書},
  pdfauthor={社内研修教材},
  pdfkeywords={情報工学, 物理, コンピュータ, プログラミング, ネットワーク, 機械学習}
]{hyperref}

% TikZ共通スタイル
\tikzset{
  box/.style={draw=RuleGray,fill=CardWhite,rounded corners=1mm,minimum height=8mm,inner sep=2mm,align=center,text=DeepNavy},
  accentbox/.style={box,draw=MutedBlue,fill=PaleBlue},
  flow/.style={-{Latex[length=2mm]},draw=MutedBlue,line width=0.5pt},
  softflow/.style={-{Latex[length=2mm]},draw=PaleGray,line width=0.5pt},
  node/.style={circle,draw=MutedBlue,fill=CardWhite,minimum size=6mm,inner sep=0pt}
}
